% Unnumbered research calculation, not a promoted theorem or a hereditary count.
\subsection*{An exact toric model for root-normal cancellation on an isolated A-chain}

Let $Y\to X$ be the minimal regular toric resolution of
$X=\operatorname{Spec}k[[x,y,z]]/(xy-z^{L+1})$, over an algebraically
closed field of characteristic $p$. Its exceptional chain has $L$
curves of square $-2$. Let $h$ be prime to $p$. Include the two
nonexceptional toric boundary branches in the boundary divisor, and
root every toric boundary component with order $h$.

The following is an exact calculation for the specified toric gerbe.
It is not an assertion identifying an arbitrary formal order along a
contracted chain with this model.

\paragraph{The finite cover and actual anchor.}
The fan may be written with rays $v_i=(1,i)$, $0\le i\le L+1$.
Replace the lattice $N$ by $hN$ and the primitive rays by $hv_i$.
The resulting regular toric resolution $Y_h$ is abstractly the same
$A_L$ resolution. Inclusion $hN\subset N$ gives a finite flat map
\[
                         q:Y_h\longrightarrow Y
\]
of degree $h^2$. On each regular toric chart it is exactly
$(u,v)\mapsto(u^h,v^h)$. Its group is $G=\mu_h^2$, and
$[Y_h/G]$ is the mixed root stack with common order $h$ along the
entire toric boundary. This description also holds after completion
at the unique point over the singular vertex.

Fix a nondegenerate scalar cocycle on $G$, with commutator a primitive
$h$th root of unity. The shift and clock matrices give a projective
representation $V_h$ of dimension $h$. The trivial bundle
$\mathcal O_{Y_h}\otimes V_h$ with this projective action is an actual
weight-one bundle $A_h$ on the corresponding coefficient gerbe.
Changing the cocycle to another primitive commutator merely permutes
the characters below. A larger scalar band can be retained by the
usual pushout of the central extension. A scalar line twist does not
alter the endomorphism calculation.

For $M=kh$, put $A=A_h^{\oplus k}$. Let $B$ be its coarse
endomorphism algebra. Conjugation on $\operatorname{End}(V_h)$ is
the regular character representation of $G$: the $h^2$ matrices
$X^aZ^b$, $0\le a,b<h$, are eigenvectors of distinct characters.
Consequently, as an \emph{ordinary module},
\[
              B\simeq\operatorname{Mat}_k(q_*\mathcal O_{Y_h}).
\]
The multiplication on $B$ is the quantum multiplication induced by
the cocycle; this displayed isomorphism is not claimed to be an
isomorphism of algebras with a commutative-cover matrix algebra.
It follows by matching each character summand of $q_*\mathcal O$
with its unique inverse character in $\operatorname{End}(V_h)$.
This is where the factor $h^2$ cancels against $k^2=M^2/h^2$ in
cycle cohomology.

\paragraph{Whole-cycle cohomology in this model.}
Let $Z>0$ be an integral cycle supported on the exceptional chain,
and denote by $Z_h$ the cycle with the same coefficients on $Y_h$.
The exact chart equations show $q^*Z=hZ_h$. Flat finite base change
and the displayed ordinary-module decomposition give
\[
 H^i(Z,B|_Z)\simeq
 \operatorname{Mat}_k H^i(hZ_h,\mathcal O_{hZ_h}).
\]
The $A_L$ singularity is rational in every characteristic. Its
regular resolution has $R^1\pi_*\mathcal O=0$, and coherent $H^2$
on the proper local resolution vanishes. Thus every effective
exceptional cycle $D$ satisfies $H^1(D,\mathcal O_D)=0$. The exact
nonreduced Riemann--Roch formula and $K_Y\cdot E_i=0$ give
$h^0(\mathcal O_D)=-D^2/2$. Applying this to $D=hZ_h$ proves
\[
       h^1(Z,B|_Z)=0,\qquad
       h^0(Z,B|_Z)=\frac{M^2}{2}(-Z^2).
\]
No term proportional to the length of the chain times the largest
cycle coefficient occurs. These cohomology identities commute with
arbitrary extension of the coefficient field to a $k$-algebra,
because all geometry and modules in this calculation are fixed.
They do not assert effectivity for an arbitrary varying order-module
family.

For completeness, if $Z$ is anti-nef on the chain and $j\in\mathbf Z$,
one also has the exact bounded-power formula
\[
 h^0(Z,B(jZ)|_Z)+h^1(Z,B(jZ)|_Z)
                 =M^2\left|j-\frac12\right|(-Z^2).
\]
Indeed, for an effective anti-nef cycle $D$ and $j\ge1$, a nonzero
section of $\mathcal O_D(jD)$ would, by its exact Cartier
annihilator $0<T\le D$, yield an injection
$\mathcal O_T\hookrightarrow\mathcal O_T((j-1)D+T)$ with
zero-dimensional cokernel. The degree would be
$(j-1)D\cdot T+T^2<0$, a contradiction. Hence $h^0=0$ for
$j\ge1$. For $j\le0$, curve duality and
$\omega_D\simeq\mathcal O_D(D)$ give $h^1=0$. Riemann--Roch
then gives the formula, and the finite cover replaces $D$ by
$hZ_h$ as above.

\paragraph{Precisely what is still missing globally.}
\begin{enumerate}
\item The actual coefficient gerbe and chosen anchor on a formal
      two-ended chain must be compared to this model, including
      their fixed boundary framings. Equality of residue orders
      alone is not used here as such an identification theorem.
\item A balanced toric anchor need not be the anchor stable under
      the canonical full order. No uniform comparison pole or
      simultaneous Frobenius bound between those choices is proved.
\item The energy above is the \emph{principal chain energy}
      $-Z_{\mathrm{chain}}^2$. Summing it over cut-out chains need
      not be bounded by the whole-cycle energy: deleting the core
      cross terms can lose precisely the required cancellation.
      A Schur-complement or finite-boundary gluing argument is
      still required.
\item The calculation bounds the fixed anchor order, not the
      generation or short projective resolutions of all locally
      free $B$-modules being counted. It supplies no $B^{s+1}$
      quotient for an arbitrary such module.
\end{enumerate}
Thus this is an exact model calculation supporting a possible
ramified A-chain contraction. It is not an unbounded-period
hereditary atlas or a main-gate reduction.

\paragraph{Source scope.}
The finite toric map and character computation are given explicitly.
Rational $A_L$ resolution cohomology, nonreduced cycle
Riemann--Roch and coherent $H^2$ vanishing are the existing fully
audited RDP/cycle inputs. The exact-annihilator argument is the
proved exceptional Picard lemma. This unnumbered note has not
received four independent audits and introduces no new primary
source or numbered statement.
